%% toc: P1 8 8.3,8.4,8.7
\section{Proofs for substitution, equality replacement and equivalence signs}
\label{sec:annex-p01-w02-p01-025-equality}

Throughout, derivations are those of \(\mathsf{ND}_{\mathsf{FOL}}\) or
\(\mathsf{ND}^{-}_{\mathsf{FOL}}\) under the conventions of
\cref{def:p01-w02-p01-025-nd-conventions,def:p01-w02-p01-025-eq-axioms}.
Two facts about term substitution from \cref{lem:substitution-fv} are used
repeatedly: if \(x\notin\operatorname{Var}(u)\) then \(u[s/x]=u\), and
\(\operatorname{Var}(u[s/x])\subseteq(\operatorname{Var}(u)\setminus\{x\})
\cup\operatorname{Var}(s)\).  Substitution in an atomic formula substitutes
in each term, so for a variable \(w\) that occurs in an atomic formula
\(\theta\) only at one argument place, \(\theta[s/w]\) is \(\theta\) with
that occurrence replaced by \(s\).

\begin{annexitem}{annex:p01-w02-p01-025-symmetry}{thm:p01-w02-p01-025-symmetry}{Symmetry of identity}
\begin{proof}
Choose a variable \(x\notin\operatorname{Var}(r)\) and put
\(\varphi\defeq(x=r)\).  Then \(\varphi[r/x]\) is \(r=r\), because
\(x[r/x]=r\) and \(r[r/x]=r\), and \(\varphi[t/x]\) is \(t=r\), because
\(r[t/x]=r\).  The derivation is
\[
\begin{array}{rll}
1.&r=t&\text{assumption}\\
2.&r=r&=I\\
3.&t=r&=E,\ 1,2.
\end{array}
\]
Its only open assumption is \(r=t\).
\end{proof}
\end{annexitem}

\begin{annexitem}{annex:p01-w02-p01-025-transitivity}{thm:p01-w02-p01-025-transitivity}{Transitivity of identity}
\begin{proof}
Choose \(x\notin\operatorname{Var}(r)\) and put \(\varphi\defeq(r=x)\).
Then \(\varphi[t/x]\) is \(r=t\) and \(\varphi[u/x]\) is \(r=u\).  The
derivation is
\[
\begin{array}{rll}
1.&t=u&\text{assumption}\\
2.&r=t&\text{assumption}\\
3.&r=u&=E,\ 1,2,
\end{array}
\]
with open assumptions \(r=t\) and \(t=u\).
\end{proof}
\end{annexitem}

\begin{annexitem}{annex:p01-w02-p01-025-fun-cong}{thm:p01-w02-p01-025-fun-cong}{Function congruence schema}
\begin{proof}
If \(n=0\), the one-line derivation \(c=c\) by \(=I\) has no open
assumptions.

Let \(n\geq 1\).  Choose a variable \(w\) that occurs in none of
\(r_1,\ldots,r_n,t_1,\ldots,t_n\).  For \(0\leq k\leq n\) put
\[
 a_k\defeq f(t_1,\ldots,t_k,r_{k+1},\ldots,r_n),
\]
so \(a_0=f(r_1,\ldots,r_n)\) and \(a_n=f(t_1,\ldots,t_n)\).  We build, for
each \(k\), a derivation \(D_k\) of \(a_0=a_k\) whose open assumptions are
exactly \(r_1=t_1,\ldots,r_k=t_k\), each occurring once as a leaf.

\emph{Step \(0\).}  \(D_0\) is the single application of \(=I\) with
conclusion \(a_0=a_0\).

\emph{Step \(k+1\), for \(k<n\).}  Put
\[
 \varphi_k\defeq
 \bigl(a_0=f(t_1,\ldots,t_k,w,r_{k+2},\ldots,r_n)\bigr).
\]
The variable \(w\) occurs in \(\varphi_k\) at exactly one place, the
\((k+1)\)-st argument on the right; it does not occur in \(a_0\) or in any
other argument.  Hence \(\varphi_k[r_{k+1}/w]\) is \(a_0=a_k\) and
\(\varphi_k[t_{k+1}/w]\) is \(a_0=a_{k+1}\).  Apply \(=E\) with major premise
the assumption \(r_{k+1}=t_{k+1}\) and minor premise the conclusion of
\(D_k\); the result is \(D_{k+1}\).  Its open assumptions are those of
\(D_k\) together with \(r_{k+1}=t_{k+1}\).

\(D_n\) is the required derivation.  It contains one application of \(=I\)
and \(n\) applications of \(=E\), and no rule with a discharge or an
eigenvariable.
\end{proof}
\end{annexitem}

\begin{annexitem}{annex:p01-w02-p01-025-rel-cong}{thm:p01-w02-p01-025-rel-cong}{Relation congruence schema}
\begin{proof}
If \(n=0\), the one-node derivation consisting of the assumption \(R\) has
conclusion \(R\).

Let \(n\geq 1\), choose \(w\) as in the function case, and put
\(b_k\defeq R(t_1,\ldots,t_k,r_{k+1},\ldots,r_n)\) for
\(0\leq k\leq n\).  \(D_0\) is the assumption \(b_0\), that is,
\(R(r_1,\ldots,r_n)\).
Given \(D_k\) with conclusion \(b_k\), put
\[
 \varphi_k\defeq R(t_1,\ldots,t_k,w,r_{k+2},\ldots,r_n).
\]
Then \(\varphi_k[r_{k+1}/w]=b_k\) and \(\varphi_k[t_{k+1}/w]=b_{k+1}\).  One
application of \(=E\) to the assumption \(r_{k+1}=t_{k+1}\) and the
conclusion of \(D_k\) gives \(D_{k+1}\).  After \(n\) steps the conclusion is
\(b_n=R(t_1,\ldots,t_n)\), and the open assumptions are
\(R(r_1,\ldots,r_n)\) and \(r_1=t_1,\ldots,r_n=t_n\).

For the identity sign the same two steps apply to the atomic formulas
\(\varphi_0\defeq(w=r_2)\) and \(\varphi_1\defeq(t_1=w)\): from
\(r_1=r_2\) and \(r_1=t_1\) the first step gives \(t_1=r_2\), and from that
and \(r_2=t_2\) the second gives \(t_1=t_2\).
\end{proof}
\end{annexitem}

\begin{annexitem}{annex:p01-w02-p01-025-term-subst}{thm:p01-w02-p01-025-term-subst}{Substitutivity in terms}
\begin{proof}
Induct on the term \(u\).  If \(u\) is \(x\), then \(u[r/x]=r\) and
\(u[t/x]=t\), and the assumption \(r=t\) is the derivation.  If \(u\) is a
variable \(y\ne x\), both sides are \(y\), and \(=I\) derives \(y=y\).  If
\(u=f(u_1,\ldots,u_n)\), then \(u[s/x]=f(u_1[s/x],\ldots,u_n[s/x])\) for
\(s\in\{r,t\}\).  By the induction hypothesis there are derivations
\(E_i\) of \(u_i[r/x]=u_i[t/x]\) with open assumptions among \(\{r=t\}\).
\Cref{thm:p01-w02-p01-025-fun-cong}, with \(r_i\defeq u_i[r/x]\) and
\(t_i\defeq u_i[t/x]\), gives a derivation \(D\) of \(u[r/x]=u[t/x]\) whose
open assumptions are the leaves \(u_i[r/x]=u_i[t/x]\).  Place \(E_i\) above
each such leaf.  Because \(D\) uses only \(=I\) and \(=E\), which have
neither discharges nor conditions on open assumptions, the result is a
derivation, and its open assumptions lie in \(\{r=t\}\).  For \(n=0\), \(D\)
is the \(=I\) step \(c=c\).
\end{proof}
\end{annexitem}

\begin{annexitem}{annex:p01-w02-p01-025-axioms-derivable}{thm:p01-w02-p01-025-axioms-derivable}{The equality rules derive the equality axioms}
\begin{proof}
\emph{Reflexivity.}  \(=I\) gives \(x=x\) with no open assumption, and
\(\forall I\) gives \(\forall x\,(x=x)\); the eigenvariable condition holds
because there is no open assumption.

\emph{Function axioms.}  For \(n=0\) the axiom \(c=c\) is an \(=I\) step.
For \(n\geq 1\), apply \cref{thm:p01-w02-p01-025-fun-cong} with
\(r_i\defeq x_i\) and \(t_i\defeq y_i\).  The resulting derivation \(D_n\)
of \(f(x_1,\ldots,x_n)=f(y_1,\ldots,y_n)\) has open assumptions exactly
\(x_1=y_1,\ldots,x_n=y_n\), each occurring once; these formulas are pairwise
distinct because the variables are distinct.  Apply \(\to I\) \(n\) times,
discharging \(x_n=y_n\), then \(x_{n-1}=y_{n-1}\), and so on down to
\(x_1=y_1\).  Each discharge marks an open occurrence, so none is vacuous.
The result has no open assumptions, and \(2n\) applications of \(\forall I\),
generalizing \(y_n,\ldots,y_1\) and then \(x_n,\ldots,x_1\), are permitted.

\emph{Relation axioms.}  For a nullary \(R\), assume \(R\) and discharge it by
\(\to I\) to obtain \(R\to R\).  For \(n\geq 1\), apply
\cref{thm:p01-w02-p01-025-rel-cong} with \(r_i\defeq x_i\) and
\(t_i\defeq y_i\), for \(R\in\mathcal R\) or for the identity sign.  The open
assumptions are \(R(x_1,\ldots,x_n)\) and the \(n\) equations, each
occurring as a leaf.  Discharge \(R(x_1,\ldots,x_n)\) first and then the
equations in the order \(x_n=y_n,\ldots,x_1=y_1\), and generalize as before.
\end{proof}
\end{annexitem}

\begin{annexitem}{annex:p01-w02-p01-025-axiom-replacement}{thm:p01-w02-p01-025-axiom-replacement}{Replacement from the equality axioms}
\begin{proof}
All derivations in this proof are in \(\mathsf{ND}^{-}_{\mathsf{FOL}}\), and
\(\derives\) abbreviates \(\derives_{\mathsf{ND}^{-}_{\mathsf{FOL}}}\).
Write \(\mathrm{Eq}\) for \(\mathrm{Eq}_{\mathcal L}\).

\emph{Substitution sequences.}  A \emph{substitution sequence} is a finite
list \(\sigma=\langle[s_1/z_1],\ldots,[s_k/z_k]\rangle\) with \(k\geq 0\).
For a term or formula \(e\), put
\(e\sigma\defeq e[s_1/z_1]\cdots[s_k/z_k]\), applied from left to right, and
write \(\sigma\cdot[s/z]\) for the extended list.  Because substitution
commutes with every connective and acts termwise on atomic formulas, an
induction on \(k\) shows that the following pairs are identical strings, for
\(\circ\in\{\land,\lor,\to\}\):
\[
\begin{array}{ll}
 R(u_1,\ldots,u_n)\sigma & R(u_1\sigma,\ldots,u_n\sigma),\\
 (u_1=u_2)\sigma & (u_1\sigma=u_2\sigma),\\
 \bottom\sigma & \bottom,\\
 (\neg\chi)\sigma & \neg(\chi\sigma),\\
 (\chi\circ\theta)\sigma & (\chi\sigma\circ\theta\sigma).
\end{array}
\]
Elsewhere in this proof an equation between formulas or terms outside an
object-language formula asserts identity of strings.

\emph{Claim A (variables of a substituted term).}  For every term \(u\),
\(\operatorname{Var}(u[s/z])=\bigcup_{v\in\operatorname{Var}(u)}
\operatorname{Var}(v[s/z])\).  For a variable this is trivial, and for
\(f(u_1,\ldots,u_n)\) both sides are the union over the arguments.

\emph{Claim B (free variables are carried).}  If \(u\in\FV(\chi)\), then
\(\operatorname{Var}(u[s/z])\subseteq\FV(\chi[s/z])\).  Strong induction on
\(c(\chi)\).  If \(\chi\) is atomic, \(u\) occurs in some term \(w\) of
\(\chi\), and \(\chi[s/z]\) contains \(w[s/z]\), whose variables include
those of \(u[s/z]\) by Claim A; the atomic \(\FV\) clause gives the claim.
\(\bottom\) has no free variable, and the connective cases follow from the
\(\FV\) clauses.  Let \(\chi=Qy\,\psi\), so \(u\ne y\) and \(u\in\FV(\psi)\).
If \(y=z\), then \(\chi[s/z]=\chi\) and \(u[s/z]=u\).  If \(y\ne z\) and
\(y\notin\operatorname{Var}(s)\), then \(\chi[s/z]=Qy\,(\psi[s/z])\); by the
induction hypothesis \(\operatorname{Var}(u[s/z])\subseteq\FV(\psi[s/z])\),
and \(y\notin\operatorname{Var}(u[s/z])\subseteq\{u\}\cup
\operatorname{Var}(s)\).  If \(y\ne z\) and \(y\in\operatorname{Var}(s)\),
then \(\chi[s/z]=Qv\,(\psi[v/y][s/z])\) with
\(v\notin\operatorname{Var}(\psi)\cup\operatorname{Var}(s)\cup\{z,y\}\).
Since \(u\ne y\), the induction hypothesis for \(\psi\) gives
\(u\in\FV(\psi[v/y])\); since \(c(\psi[v/y])=c(\psi)\) by
\cref{lem:substitution-complexity}, the induction hypothesis also applies to
\(\psi[v/y]\) and gives
\(\operatorname{Var}(u[s/z])\subseteq\FV(\psi[v/y][s/z])\).  Finally
\(v\notin\operatorname{Var}(u[s/z])\subseteq\{u\}\cup\operatorname{Var}(s)\),
because \(v\ne u\in\operatorname{Var}(\psi)\).  Iterating Claims A and B
along a sequence gives: if \(u\in\FV(\chi)\), then
\(\operatorname{Var}(u\sigma)\subseteq\FV(\chi\sigma)\).

\emph{Claim C (a sequence applied to a quantified formula).}  For every
formula \(Qy\,\psi\) and sequence \(\sigma\) there are a variable
\(y_\sigma\) and a sequence \(\sigma'\) such that
\((Qy\,\psi)\sigma=Qy_\sigma\,(\psi\sigma')\) and
(a) \(y\sigma'=y_\sigma\); (b) \(u\sigma'=u\sigma\) for every
\(u\in\FV(Qy\,\psi)\); (c) \(y_\sigma\notin\operatorname{Var}(u\sigma)\) for
every \(u\in\FV(Qy\,\psi)\).  Induct on the length of \(\sigma\).  For the
empty sequence take \(y_\sigma=y\) and \(\sigma'\) empty; (c) holds because
\(u\ne y\).  Let \(\sigma=\rho\cdot[s/x]\), with
\((Qy\,\psi)\rho=Qy_\rho\,\chi\), \(\chi=\psi\rho'\), and (a)--(c) for
\(\rho\).  Apply the three clauses of the substitution definition to
\(Qy_\rho\,\chi\).
\begin{itemize}
 \item \(y_\rho=x\): the result is \(Qy_\rho\,\chi\).  Take
 \(y_\sigma=y_\rho\), \(\sigma'=\rho'\).  By (c) for \(\rho\),
 \(x\notin\operatorname{Var}(u\rho)\), so \(u\sigma=u\rho[s/x]=u\rho\), and
 (a)--(c) carry over.
 \item \(y_\rho\ne x\), \(y_\rho\notin\operatorname{Var}(s)\): the result is
 \(Qy_\rho\,(\chi[s/x])\).  Take \(y_\sigma=y_\rho\),
 \(\sigma'=\rho'\cdot[s/x]\).  Then \(y\sigma'=y_\rho[s/x]=y_\rho\);
 \(u\sigma'=u\rho[s/x]=u\sigma\); and
 \(y_\rho\notin\operatorname{Var}(u\rho[s/x])\subseteq
 \operatorname{Var}(u\rho)\cup\operatorname{Var}(s)\).
 \item \(y_\rho\ne x\), \(y_\rho\in\operatorname{Var}(s)\): the result is
 \(Qv\,(\chi[v/y_\rho][s/x])\) with \(v\notin\operatorname{Var}(\chi)\cup
 \operatorname{Var}(s)\cup\{x,y_\rho\}\).  Take \(y_\sigma=v\) and
 \(\sigma'=\rho'\cdot[v/y_\rho]\cdot[s/x]\).  Then
 \(y\sigma'=y_\rho[v/y_\rho][s/x]=v[s/x]=v\).  For \(u\in\FV(Qy\,\psi)\),
 \(u\rho'=u\rho\) does not contain \(y_\rho\), so
 \(u\sigma'=u\rho[s/x]=u\sigma\).  By Claim B,
 \(\operatorname{Var}(u\rho)=\operatorname{Var}(u\rho')\subseteq
 \FV(\psi\rho')\subseteq\operatorname{Var}(\chi)\), so
 \(v\notin\operatorname{Var}(u\rho)\cup\operatorname{Var}(s)\supseteq
 \operatorname{Var}(u\sigma)\).
\end{itemize}

\emph{Claim D (instantiation).}  Let \(\theta\) be quantifier-free, let
\(v_1,\ldots,v_m\) be distinct variables, and let \(s_1,\ldots,s_m\) be terms
in which no \(v_j\) occurs.  Then \(m\) applications of \(\forall E\), with
the terms \(s_1,\ldots,s_m\) in order, derive from
\(\forall v_1\cdots\forall v_m\,\theta\) the formula obtained from \(\theta\)
by replacing each occurrence of each \(v_j\) by \(s_j\).  At the \(j\)-th step
every remaining binder \(v_l\) (\(l>j\)) differs from \(v_j\) and does not
occur in \(s_j\), so the second clause of the substitution definition applies
at each quantifier, and on the quantifier-free matrix substitution acts
termwise; terms inserted earlier contain no \(v_j\).

\emph{Claim E (discharge and composition).}  (i) If
\(\Sigma\cup\{\alpha\}\derives\beta\), then
\(\Sigma\derives\alpha\to\beta\).  If the given derivation has an open
occurrence of \(\alpha\), discharge all of them by \(\to I\).  Otherwise it
derives \(\beta\) from \(\Sigma\); assume \(\alpha\), form
\(\beta\land\alpha\) by \(\land I\), recover \(\beta\) by \(\land E_1\), and
discharge the new occurrence of \(\alpha\).  (ii) If
\(\Sigma\cup\{\alpha\}\derives\beta\) and \(\Sigma'\derives\alpha\), then
\(\Sigma\cup\Sigma'\derives\beta\), by (i) and \(\to E\).  Neither step has
an eigenvariable.

\emph{Claim F (renaming the axioms used).}  If
\(\Sigma\cup E_0\derives\beta\) with \(E_0\subseteq\mathrm{Eq}\) finite and
\(V\) is a finite set of variables, then \(\Sigma\cup E_0'\derives\beta\) for
a finite \(E_0'\subseteq\mathrm{Eq}\) none of whose members contains a
variable of \(V\).  For \(\varepsilon\in E_0\) with prefix variables
\(v_1,\ldots,v_m\), choose distinct \(v_1',\ldots,v_m'\) outside
\(V\cup\operatorname{Var}(\varepsilon)\), and let \(\varepsilon'\in
\mathrm{Eq}\) be the axiom of the same kind written with them.  Renaming
\(v_1\) to \(v_1'\), then \(v_2\) to \(v_2'\), and so on, is a chain of
generating steps of \(\equiv_\alpha\) applied to a subformula, licensed
because each new variable is absent from that subformula, and lifted by the
congruence clause; so \(\varepsilon'\equiv_\alpha\varepsilon\), and the
alpha rule derives \(\varepsilon\) from \(\varepsilon'\).  Claim E(ii)
replaces each \(\varepsilon\) by \(\varepsilon'\).

\emph{Claim G (symmetry and reflexivity from the axioms).}  For all terms
\(a,b\): \(\mathrm{Eq}\derives a=a\) and
\(\mathrm{Eq}\cup\{a=b\}\derives b=a\).  The first is Claim D applied to a
reflexivity axiom whose bound variable does not occur in \(a\).  For the
second take the relation axiom for \(=\) with bound variables outside
\(\operatorname{Var}(a)\cup\operatorname{Var}(b)\) and instantiate
\(x_1,x_2,y_1,y_2\) by \(a,a,b,a\) (Claim D).  This gives
\(a=b\to(a=a\to(a=a\to b=a))\), and three applications of \(\to E\), with
\(a=b\) and twice \(a=a\), give \(b=a\).

\emph{Claim H (terms).}  Let \(H\) be a finite set of formulas and
\(\sigma,\tau\) sequences.  If \(\mathrm{Eq}\cup H\derives v\sigma=v\tau\)
for every \(v\in\operatorname{Var}(u)\), then
\(\mathrm{Eq}\cup H\derives u\sigma=u\tau\).  Induct on \(u\).  A variable is
the hypothesis.  For \(u=f(u_1,\ldots,u_n)\),
\(u\sigma=f(u_1\sigma,\ldots,u_n\sigma)\), and likewise for \(\tau\).  Take
the function axiom for \(f\) with bound variables outside every
\(\operatorname{Var}(u_i\sigma)\cup\operatorname{Var}(u_i\tau)\), instantiate
it by Claim D to
\(u_1\sigma=u_1\tau\to(\cdots\to f(u_1\sigma,\ldots)=f(u_1\tau,\ldots))\),
and apply \(\to E\) \(n\) times to the derivations given by the induction
hypothesis.  For \(n=0\) the axiom \(c=c\) is the derivation.

\emph{Main claim.}  Say that \((\sigma,\tau)\) is \emph{\(H\)-matched on}
\(\varphi\) if \(\mathrm{Eq}\cup H\derives u\sigma=u\tau\) for every
\(u\in\FV(\varphi)\).  We prove by strong induction on \(c(\varphi)\), for
all finite \(H\) and all sequences \(\sigma,\tau\): if \((\sigma,\tau)\) is
\(H\)-matched on \(\varphi\), then
\[
 \mathrm{Eq}\cup H\cup\{\varphi\sigma\}\derives\varphi\tau.
\]
By Claim G and Claim E(ii), a pair matched on \(\varphi\) is also matched in
the order \((\tau,\sigma)\), and matched on every subformula with no more
free variables.

\emph{Atomic.}  Let \(\varphi=R(u_1,\ldots,u_n)\), with \(R\in\mathcal R\) or
\(R\) the identity sign.  Every variable of each \(u_i\) is free in
\(\varphi\), so Claim H gives \(\mathrm{Eq}\cup H\derives
u_i\sigma=u_i\tau\).  Instantiate the relation axiom for \(R\), with bound
variables outside the variables of all \(u_i\sigma,u_i\tau\), by Claim D, and
apply \(\to E\) \(n\) times to these derivations and once to the assumption
\(\varphi\sigma=R(u_1\sigma,\ldots,u_n\sigma)\).  The conclusion is
\(R(u_1\tau,\ldots,u_n\tau)=\varphi\tau\).  For nullary \(R\) the axiom
\(R\to R\) and the assumption \(R\) suffice.  For \(\bottom\),
\(\bottom\sigma=\bottom=\bottom\tau\) and the assumption is the derivation.

\emph{Negation.}  \(\varphi\sigma=\neg(\chi\sigma)\).  The induction
hypothesis in the order \((\tau,\sigma)\) gives
\(\mathrm{Eq}\cup H\cup\{\chi\tau\}\derives\chi\sigma\).  With the assumption
\(\neg(\chi\sigma)\), \(\neg E\) gives \(\bot\).  If \(\chi\tau\) is open in
this derivation, \(\neg I\) discharges it and yields \(\neg(\chi\tau)\);
otherwise \(\bot E\) yields \(\neg(\chi\tau)\) directly.

\emph{Conjunction.}  \(\land E_1\) and \(\land E_2\) give \(\chi\sigma\) and
\(\theta\sigma\); the induction hypotheses and Claim E(ii) give \(\chi\tau\)
and \(\theta\tau\); \(\land I\) concludes.

\emph{Disjunction.}  From \(\chi\sigma\) the induction hypothesis and
\(\lor I_1\) give \(\chi\tau\lor\theta\tau\), and from \(\theta\sigma\) the
induction hypothesis and \(\lor I_2\) give the same formula.  If both case
assumptions are open in their subderivations, \(\lor E\) applied to the
assumption \(\varphi\sigma\) discharges them.  If one is not open, that
subderivation already derives \(\varphi\tau\) from \(\mathrm{Eq}\cup H\).

\emph{Implication.}  Assume \(\chi\tau\).  The induction hypothesis in the
order \((\tau,\sigma)\) gives \(\chi\sigma\); \(\to E\) with
\(\varphi\sigma=(\chi\sigma\to\theta\sigma)\) gives \(\theta\sigma\); the
induction hypothesis gives \(\theta\tau\).  Claim E(i) discharges
\(\chi\tau\) and yields \(\varphi\tau\).

\emph{Quantifiers.}  Let \(\varphi=Qy\,\psi\).  By Claim C,
\(\varphi\sigma=Qy_\sigma\,(\psi\sigma')\) and
\(\varphi\tau=Qy_\tau\,(\psi\tau')\) with (a)--(c).  Choose a variable \(v\)
that occurs in none of \(\varphi\sigma\), \(\varphi\tau\) and the members of
\(H\).  Put \(\sigma''\defeq\sigma'\cdot[v/y_\sigma]\) and
\(\tau''\defeq\tau'\cdot[v/y_\tau]\), so
\(\psi\sigma''=(\psi\sigma')[v/y_\sigma]\) and
\(\psi\tau''=(\psi\tau')[v/y_\tau]\).  The pair \((\sigma'',\tau'')\) is
\(H\)-matched on \(\psi\): by (a), \(y\sigma''=y\tau''=v\), and
\(\mathrm{Eq}\derives v=v\) by Claim G; for \(u\in\FV(\psi)\) with
\(u\ne y\), (b) and (c) give \(u\sigma''=u\sigma[v/y_\sigma]=u\sigma\) and
likewise \(u\tau''=u\tau\), and \(u\in\FV(\varphi)\).  Since
\(c(\psi)<c(\varphi)\), the induction hypothesis gives a derivation \(D\) of
\(\psi\tau''\) from \(\mathrm{Eq}\cup H\cup\{\psi\sigma''\}\).  By Claim F we
may assume that no axiom open in \(D\) contains \(v\).

For \(Q=\forall\): \(\forall E\) with the term \(v\) derives
\((\psi\sigma')[v/y_\sigma]=\psi\sigma''\) from \(\varphi\sigma\); Claim E(ii)
gives \(\psi\tau''\) from \(\mathrm{Eq}\cup H\cup\{\varphi\sigma\}\).  The
variable \(v\) is free in no open assumption: the axioms are sentences, and
\(v\) occurs neither in \(H\) nor in \(\varphi\sigma\).  So \(\forall I\)
gives \(\forall v\,((\psi\tau')[v/y_\tau])\).  Since
\(v\notin\operatorname{Var}(\psi\tau')\), the generating clause gives
\(Qy_\tau\,(\psi\tau')\equiv_\alpha Qv\,((\psi\tau')[v/y_\tau])\), and the
alpha rule yields \(\varphi\tau\).

For \(Q=\exists\): \(\exists I\) applied to the conclusion
\((\psi\tau')[v/y_\tau]\) of \(D\) yields \(\exists y_\tau\,(\psi\tau')=
\varphi\tau\).  If \(\psi\sigma''\) is not open in \(D\), this derives
\(\varphi\tau\) from \(\mathrm{Eq}\cup H\).  Otherwise apply \(\exists E\)
with major premise the assumption
\(\varphi\sigma=\exists y_\sigma\,(\psi\sigma')\), minor assumption
\((\psi\sigma')[v/y_\sigma]=\psi\sigma''\) and eigenvariable \(v\).  The
side condition holds: \(v\) is absent from the conclusion \(\varphi\tau\),
from \(\varphi\sigma\), from \(H\), and from every axiom open in \(D\).

\emph{Conclusion.}  Take \(H=\{r=t\}\), \(\sigma=\langle[r/x]\rangle\) and
\(\tau=\langle[t/x]\rangle\).  For \(u=x\), \(u\sigma=r\) and \(u\tau=t\),
and \(r=t\in H\); for a free variable \(u\ne x\), \(u\sigma=u=u\tau\), and
Claim G gives \(u=u\).  The pair is therefore \(H\)-matched on \(\varphi\),
and the main claim gives
\(\mathrm{Eq}\cup\{r=t,\varphi[r/x]\}\derives\varphi[t/x]\).
\end{proof}
\end{annexitem}

\begin{annexitem}{annex:p01-w02-p01-025-axiom-metatheorem}{thm:p01-w02-p01-025-axiom-metatheorem}{Equality rules and equality axioms have the same strength}
\begin{proof}
\emph{Right to left.}  Every rule instance of
\(\mathsf{ND}^{-}_{\mathsf{FOL}}\) is a rule instance of
\(\mathsf{ND}_{\mathsf{FOL}}\) with the same side conditions, so a derivation
of \(\varphi\) from \(\Gamma_0\cup E_0\), with \(\Gamma_0\subseteq\Gamma\) and
\(E_0\subseteq\mathrm{Eq}_{\mathcal L}\) finite, is an
\(\mathsf{ND}_{\mathsf{FOL}}\)-derivation.  Each \(\varepsilon\in E_0\) has
an \(\mathsf{ND}_{\mathsf{FOL}}\)-derivation without open assumptions by
\cref{thm:p01-w02-p01-025-axioms-derivable}.  Claim E(ii) of
Annex~\ref{annex:p01-w02-p01-025-axiom-replacement}, whose proof uses only
\(\to I\), \(\to E\), \(\land I\) and \(\land E_1\) and so holds in
\(\mathsf{ND}_{\mathsf{FOL}}\) as well, removes the members of \(E_0\) one at a
time.

\emph{Left to right.}  By induction on an
\(\mathsf{ND}_{\mathsf{FOL}}\)-derivation \(D\) we construct an
\(\mathsf{ND}^{-}_{\mathsf{FOL}}\)-derivation \(D^{*}\) of the same conclusion
whose open assumptions are open assumptions of \(D\) or members of
\(\mathrm{Eq}_{\mathcal L}\).  Claims E and F below refer to that annex item.
\begin{itemize}
 \item An assumption is unchanged.
 \item An \(=I\) step with conclusion \(t=t\) becomes a reflexivity axiom
 \(\forall z\,(z=z)\) followed by \(\forall E\) with the term \(t\); since
 \(z=z\) has no binder, its instance is \(t=t\).
 \item An \(=E\) step with premises \(r=t\) and \(\varphi[r/x]\) becomes the
 derivation of \cref{thm:p01-w02-p01-025-axiom-replacement}, composed by
 Claim E(ii) with the translations of the two premise derivations.
 \item The rules \(\land I\), \(\land E_1\), \(\land E_2\), \(\lor I_1\),
 \(\lor I_2\), \(\to E\), \(\neg E\), \(\bot E\), \(\forall E\), \(\exists I\)
 and \(\alpha\) have no discharge and no eigenvariable; apply the same rule to
 the translated premises.
 \item \(\to I\): if the discharged formula is open in the translated
 premise, apply \(\to I\); otherwise use Claim E(i).  \(\neg I\) and RAA: if
 the discharged formula is open, apply the rule; otherwise the translated
 premise derives \(\bot\) from the remaining assumptions, and \(\bot E\)
 gives the conclusion.  \(\lor E\): if a case assumption is not open in its
 translated subderivation, that subderivation already derives the conclusion
 from assumptions of the conclusion; otherwise apply \(\lor E\).
 \item \(\forall I\): the eigenvariable is free in no open assumption of
 \(D\) above the step, and the added axioms are sentences.
 \item \(\exists E\) with eigenvariable \(y\): first apply Claim F with
 \(V=\{y\}\) to both translated premise derivations, so that no open axiom
 contains \(y\).  If the witness assumption is open in the translated minor
 derivation, apply \(\exists E\); its side condition holds because it held
 in \(D\) and the remaining open assumptions are axioms not containing
 \(y\).  Otherwise the translated minor derivation already derives the
 conclusion.
\end{itemize}
A derivation of \(\varphi\) from finitely many members of \(\Gamma\)
therefore becomes a derivation from those members and finitely many equality
axioms.
\end{proof}
\end{annexitem}

\begin{annexitem}{annex:p01-w02-p01-025-lk-symmetric}{thm:p01-w02-p01-025-lk-symmetric}{Symmetric replacement in the sequent calculus with equality}
\begin{proof}
Let \(\Pi\) derive \(\Gamma,\varphi[t/x]\Rightarrow\Delta\).
\begin{enumerate}
 \item Apply \(=L\) to \(\Pi\) with the terms exchanged, that is, with
 \(t\) in the role of the rule's \(r\) and \(r\) in the role of its \(t\).
 The premise \(\Gamma,\varphi[t/x]\Rightarrow\Delta\) has the required
 form, and the conclusion is \(\Gamma,t=r,\varphi[r/x]\Rightarrow\Delta\).
 \item Derive \(r=t\Rightarrow t=r\) as in
 Annex~\ref{annex:p7.eqlk.replacement}: choose \(z\) occurring in neither
 \(r\) nor \(t\); apply \(=L\) with the formula \(t=z\) to the initial
 sequent \(t=r\Rightarrow t=r\), whose antecedent is \((t=z)[r/z]\), to
 obtain \(r=t,\ t=t\Rightarrow t=r\), since \((t=z)[t/z]\) is \(t=t\); then
 Cut against the initial sequent \(\Rightarrow t=t\).
 \item Cut \(r=t\Rightarrow t=r\) against the antecedent occurrence of
 \(t=r\) in the sequent of step~1.  The Cut rule with left premise
 \(r=t\Rightarrow t=r\) and right premise
 \(\Gamma,\varphi[r/x],t=r\Rightarrow\Delta\) concludes
 \(r=t,\Gamma,\varphi[r/x]\Rightarrow\Delta\), which is the required sequent
 because both sides are multisets.
\end{enumerate}
Step~1 is the one application of \(=L\) to the given derivation, step~3 is
the Cut against the fixed derivation of \(r=t\Rightarrow t=r\), and step~2
is that fixed derivation, with one \(=L\) and one Cut.
\end{proof}
\end{annexitem}

\begin{annexitem}{annex:p01-w02-p01-025-equiv-bridges}{thm:p01-w02-p01-025-equiv-bridges}{Valid passages between the equivalence notions}
\begin{proof}
(1) By the satisfaction clauses of \cref{sec:fol-semantics-02},
\(\mathcal M,s\entails(\varphi\to\psi)\land(\psi\to\varphi)\) iff
(\(\mathcal M,s\not\entails\varphi\) or \(\mathcal M,s\entails\psi\)) and
(\(\mathcal M,s\not\entails\psi\) or \(\mathcal M,s\entails\varphi\)), that
is, iff \(\mathcal M,s\entails\varphi\) holds exactly when
\(\mathcal M,s\entails\psi\) holds.  Validity requires this at every structure and assignment, which is
the definition of \(\varphi\equiv\psi\).

(2) \Cref{cor:alpha-invariance} gives \(\varphi\equiv\psi\).  The relation
\(c(\varphi)=c(\psi)\) is an equivalence relation; it is a congruence for
every constructor, because \(c\) of a compound formula is computed from
\(c\) of its immediate subformulas; and it contains every generating pair
\(Qx\,\varphi,\ Qy\,(\varphi[y/x])\), since
\(c(\varphi[y/x])=c(\varphi)\) by \cref{lem:substitution-complexity}.  By
the leastness clause, \(\equiv_\alpha\) is contained in it.  For the
counterexample, \(\mathcal M,s\entails P(x)\land P(x)\) iff
\(\mathcal M,s\entails P(x)\) by the conjunction clause, so the formulas are
logically equivalent, while \(c(P(x)\land P(x))=1\ne 0=c(P(x))\).

(3) \(\mathcal M\entails\sigma\) iff \(\mathcal M\entails\tau\), because
\(\sigma\equiv\tau\) gives agreement at every assignment.
\(\mathcal M\entails\tau\) iff \(\mathcal N\entails\tau\), because
elementarily equivalent structures satisfy the same sentences.

(4) By the definition of the defined biconditional in
\cref{sec:prop-syntax-semantics-01},
\(\alpha\leftrightarrow\beta\) is \((\alpha\to\beta)\land(\beta\to\alpha)\).
For a valuation \(v\), the truth clauses give
\(\widehat v((\alpha\to\beta)\land(\beta\to\alpha))=\mathbf T\) iff
\(\widehat v(\alpha)=\mathbf T\) implies \(\widehat v(\beta)=\mathbf T\) and
conversely, that is, iff \(\widehat v(\alpha)=\widehat v(\beta)\) by
bivalence.  Truth under every valuation is therefore equality of truth
values under every valuation, the registered CPL relation.

(5) Division with remainder gives natural numbers \(q,r\) with
\(t=nq+r\) and \(0\leq r<n\); \(r\) is the remainder of \(t\) on division
by \(n\).  By hypothesis \(n\) divides \(t-a\), and \(n\) divides \(nq\),
so \(n\) divides \((t-a)-nq=r-a\).  Since \(0\leq r<n\) and
\(0\leq a<n\), the integer \(r-a\) satisfies \(-n<r-a<n\).  The only
multiple of \(n\) strictly between \(-n\) and \(n\) is \(0\), because every
nonzero multiple \(kn\) with \(k\neq0\) has absolute value \(|k|n\geq n\).
Hence \(r-a=0\), that is, \(r=a\).
\end{proof}
\end{annexitem}
